\documentclass[10pt,a4paper]{amsart}
\usepackage[margin=19mm]{geometry}
\usepackage{amsmath,amssymb,mathtools}
\usepackage[scr=rsfs,cal=euler]{mathalfa}
\usepackage{xfrac}
\usepackage[unicode,hidelinks,bookmarks=false]{hyperref}
\newcommand{\ie}{i.e.\ }
\newcommand{\cf}{cf.\ }
\newcommand{\resp}{respectively\ }
\newcommand{\Subsection}{\subsection}
\providecommand{\nobreakdash}{\nobreak\hbox{-}}
\setlength{\emergencystretch}{2em}
\multlinegap0pt
\usepackage{bm}
\usepackage{bbm}
\usepackage{dsfont}
\usepackage{xspace}
\usepackage{cancel}
\usepackage{mathtools}
\usepackage{amsthm}
\makeatletter
\@ifundefined{theorem}{\newtheorem{theorem}{Theorem}}{}
\@ifundefined{conjecture}{\newtheorem{conjecture}[theorem]{Conjecture}}{}
\@ifundefined{definition}{\newtheorem{definition}[theorem]{Definition}}{}
\@ifundefined{proposition}{\newtheorem{proposition}[theorem]{Proposition}}{}
\@ifundefined{question}{\newtheorem{question}[theorem]{Question}}{}
\makeatother
\DeclareMathOperator*{\Res}{Res}
\title[Volumes of moduli spaces of bordered Klein surfaces]{Volumes of moduli spaces of bordered Klein surfaces}
\author{Elba Garcia-Failde, Paolo Gregori, Kento Osuga}
\date{Source: arXiv:2511.21986v2 (CC BY 4.0)}
\begin{document}
\maketitle

\noindent\emph{Source and licence.} Volumes of moduli spaces of bordered Klein surfaces. Version arXiv:2511.21986v2 (CC BY 4.0), distributed under the Creative Commons Attribution 4.0 International licence, as recorded in the arXiv licence metadata for this exact version and shown on the abstract page https://arxiv.org/abs/2511.21986v2. Full rights notice: licenses/arxiv-2511.21986-CC-BY-4.0.txt.\par\smallskip

\noindent\textit{Editorial scope.} Counted unit: the Theorem whose source label is \texttt{thm:1/2,2} in the article (source0.tex lines 1916--1920, source label \texttt{thm:1/2,2}), delivered together with the article's own complete original proof of it (source0.tex lines 1921--1951, one \texttt{proof} environment of 31 lines). No mathematics is written, repaired or re-derived: statement and proof are the article's own text line for line. Required context retained uncounted: def:RTR (lines 1604--1615); conj:RTR (lines 1622--1630); def:termwise (lines 1662--1668); conj:main (lines 1676--1684); prop:1/2,2 (lines 1882--1888). Every definition, display and label of those lines is retained unchanged. Omitted from this package and not redistributed: 1--1603, 1616--1621, 1631--1661, 1669--1675, 1685--1881, 1889--1915, 1952--2336. Nothing inside a retained excerpt is omitted or altered: every retained body is delivered exactly as the article's lines are written, so no omission receipt of any kind accompanies this package. Citations: the retained excerpts carry no citation command at all, so no bibliography entry of the article is transcribed and no reference is redistributed. Reference adaptation: none was needed and none was made; every reference inside the delivered unit resolves inside it, so no unresolved reference or citation is produced. Printed slips kept literal: no printed word, symbol, hypothesis or number of the counted statement or of its proof was altered. Layout and class: the article's own class is \texttt{amsart} with packages including geometry, fontenc, inputenc, amsmath, amssymb, amsthm, mathtools, enumerate, extarrows, stmaryrd, mathrsfs, textcomp, libertine, xcolor; the wrapper is the lane's pinned \texttt{amsart} editorial wrapper at 10pt on A4, which restates the theorem-like environments the retained text needs and adds the article's own macro definitions that the retained text needs (\texttt{\textbackslash Res}), with extra wrapper packages (bm, bbm, dsfont, xspace, cancel, mathtools). The delivered numbering is the unit's own. Assets: no retained span contains a figure environment, a graphics-inclusion command, a PDF-page inclusion, a table or a TikZ picture; no image file is copied into this package, the package declares no asset dependency and no image file is redistributed with it. Licence: version arXiv:2511.21986v2 of this article is distributed under the Creative Commons Attribution 4.0 International licence, as recorded in the arXiv licence metadata for this exact version and shown on the abstract page https://arxiv.org/abs/2511.21986v2. A full rights notice is delivered at licenses/arxiv-2511.21986-CC-BY-4.0.txt. Characters: every retained excerpt is pure ASCII, so the delivered engine, XeLaTeX with its default Latin Modern text font, reports no missing character.
\begin{definition}\label{def:RTR}
Given $\mathcal{S}^{\epsilon}_{\textup{total}}$, the \emph{refined topological recursion} recursively constructs the stable correlators $\omega_{g,n+1}$, which are meromorphic sections of $K_\Sigma^{\boxtimes (n+1)}$, for $g\in\frac12\mathbb{Z}_{\geq0}$ and $n\in\mathbb{Z}_{\geq0}$ satisfying $2g-2+n+1\geq1$, via the formula
\begin{align}
\omega_{g,n+1}(z_0,z_{[n]})\coloneqq&\left(\sum_{k\geq1}\Res_{z=\frac{k}{2}}+\sum_{i=0}^n\Res_{z=z_i}-\Res_{z=0}-\sum_{k\geq1}\Res_{z=-\frac{k}{2}}-\sum_{i=0}^n\Res_{z=-z_i}\right)\frac{\eta^z(z_1)}{4\omega_{0,1}(z)}{\rm Rec}_{g,n+1}(z,z_{[n]}),\label{RTR}
\end{align}
where $\eta^z(z_1)\coloneqq\int_{-z}^z\omega_{0,2}(z_1,\cdot)$, and
\begin{align}
{\rm Rec}_{g,n+1}(z,z_{[n]})=&\sum_{\substack{g_1+g_2=g\\J_1\sqcup J_2=[n]}}^*\omega_{g_1,1+|J_1|}(z,z_{J_1})\,\omega_{g_2,1+|J_2|}(z,z_{J_2})+\sum_{i=1}^n\frac{dx(z)dx(z_i)}{(x(z)-x(z_i))^2}\omega_{g,n}(z,z_{[n]\backslash\{i\}})\nonumber\\
&+\omega_{g-1,n+2}(z,z,z_{[n]})+\mathfrak{b}\,dx(z)d_z\frac{\omega_{g-\frac12,n+1}(z,z_{[n]})}{dx(z)}.\label{Rec}
\end{align}
Here, the symbol $*$ in the summation indicates that terms involving $\omega_{0,1}$ are omitted, whereas those involving $\omega_{0,2}$ and $\omega_{\frac12,1}$ are retained. We also write $d_z$ for the exterior derivative with respect to $z$.
\end{definition}

\begin{conjecture}\label{conj:RTR}
   For every $g\in\frac12\mathbb{Z}_{\geq0}$ and $n\in\mathbb{Z}_{\geq1}$ satisfying $2g-2+n\geq1$, let $\omega_{g,n}$ be the multidifferential defined in {\rm Definition \ref{def:RTR}}. Then $\omega_{g,n}$ satisfies the following properties:
    \begin{description}
     \item[{\rm RTR0}] \label{item:RTR0} the termwise inverse Laplace transform \textup{(Definition \ref{def:termwise})} of $\omega_{g,n}$ is convergent,
       \item[{\rm RTR1}] \label{item:RTR1}$\omega_{g,n}$ is a symmetric multidifferential viewed as a (possibly divergent) series,
       \item[{\rm RTR2}] \label{item:RTR2}$\omega_{g,n}$ has no residues with respect to any variable,
       \item[{\rm RTR3}] \label{item:RTR3}$\omega_{g,n}$ can have poles only at $z_i=-z_j$ and at $z_i=-\frac{k}{2}$ for all $i,j\in[n]$ and $k\in\mathbb{Z}_{\geq1}$.
   \end{description}
\end{conjecture}

\begin{definition}\label{def:termwise}
    Let $\omega$ be a series $\omega(z)=\sum_{k\geq0} \omega_{k}(z)$, where each $\omega_{k}(z)$ is a meromorphic differential on $\mathbb{C}$. For $L\in\mathbb{R}_{\geq0}$, we define the \emph{termwise inverse Laplace transform} $\hat{\mathcal{L}}^{-1}_L$ by
    \begin{equation*}
    \hat{\mathcal{L}}^{-1}_L.\;\omega\coloneqq\sum_{k\geq0} \mathcal{L}^{-1}_L.\,\omega_{k},
\end{equation*}
where $\mathcal{L}^{-1}_L$ denotes the usual inverse Laplace transform.
\end{definition}

\begin{question}\label{conj:main}
   Let $g\in\mathbb{Z}_{\geq0}$ and $n\in\mathbb{Z}_{\geq1}$ satisfy $2g-2+n>0$, and assume that Conjecture~\ref{conj:RTR} holds. Let $\omega_{g,n}$ denote the multidifferential produced by refined topological recursion on the refined spectral curve $\mathcal{S}^{\epsilon}_{\textup{total}}$, and let $V_{g,n}^\epsilon(L_{[n]})$ denote the total volume of the regularised moduli space of $n$-bordered Klein surfaces of genus $g$ with boundary lengths $L_{[n]}$.
  %Since $V_{g,n}^\epsilon(L_{[n]})$ is symmetric with respect to the entries of $L_{[n]}$, we assume without loss of generality that $L_1\geq L_2\geq\cdots\geq L_n$.
 Introducing a new parameter $b$ by $\mathfrak{b} = -\frac{b}{\sqrt{1+b}}$, does there exist a relation of the following form between $V_{g,n}^\epsilon(L_{[n]})$ and $\omega_{g,n}$ when $b=1$?
    \begin{equation}
     \prod_{i=1}^nL_i\cdot V_{g,n}^\epsilon(L_{[n]})= (1+b)^{g}(\hat{\mathcal{L}}_{L_1}^{-1}\cdots \hat{\mathcal{L}}_{L_n}^{-1}).\;\omega_{g,n}\Big|_{b=1}.\label{V and omega relation}
    \end{equation}
    %after appropriate analytic continuation on the right-hand side with respect to $L_{[n]}$.
\end{question}

\begin{proposition}\label{prop:1/2,2}
Conjecture~\ref{conj:RTR} holds for $\omega_{\frac12,2}$. Furthermore, it can be explicitly calculated as
\begin{align}
\omega_{\frac12,2}(z_1,z_2)=&2\mathfrak{b}\log\left(2\sinh\frac{\epsilon}{2}\right)\frac{dz_1dz_2}{z_1^2z_2^2}-\mathfrak{b}d_{z_1}\frac{\Delta_{z_1}\omega_{0,2}(z_1,z_2)}{4\omega_{0,1}(z_1)}-\mathfrak{b}d_{z_2}\frac{\Delta_{z_2}\omega_{0,2}(z_2,z_1)}{4\omega_{0,1}(z_2)}\nonumber\\
&+\frac{\mathfrak{b}}{2}\sum_{k\geq1}\frac{(-1)^k}{k}\left(\frac{dz_1}{(z_1-\frac{k}{2})^2}+\frac{dz_1}{(z_1+\frac{k}{2})^2}\right)\left(\frac{dz_2}{(z_2-\frac{k}{2})^2}+\frac{dz_2}{(z_2+\frac{k}{2})^2}\right)\label{explicit w_{1/2,2}}.
\end{align}
\end{proposition}

\begin{theorem}\label{thm:1/2,2} Question~\ref{conj:main} is answered affirmatively for $(g,n)=\big(\frac12,2\big)$. Explicitly,
 \begin{equation}
     L_1L_2\cdot V_{\frac12,2}^\epsilon(L_{[2]})= (1+b)^{\frac12}(\hat{\mathcal{L}}_{L_1}^{-1} \hat{\mathcal{L}}_{L_2}^{-1}).\; \omega_{\frac12,2}\Big|_{b=1}.
    \end{equation}
\end{theorem}

\begin{proof}
Let us assume $L_1\geq L_2$, and first consider $\mathcal{L}^{-1}_{L_2}.\;\omega_{\frac12,2}(z_1,z_2)$. Note that in general we have
\[\Res_{z=-u}\frac{1}{(z+u)^m}e^{z L} = \frac{1}{(m-1)!} L^{m-1} e^{-u L}.\]
Thus, thanks to Proposition \ref{prop:1/2,2}, the termwise inverse Laplace transform reduces to taking residues at the poles of $\omega_{\frac12,2}$. This can be explicitly computed as follows:
    \begin{align}
        \Res_{z_2=0}\omega_{\frac12,2}(z_1,z_2)e^{z_2L_2}=&2L_2\mathfrak{b}\log\left(2\sinh\frac{\epsilon}{2}\right)\frac{dz_1}{z_1^2}-\frac{\mathfrak{b}L_2^2}{2}\frac{dz_1}{z_1^2},\label{inverse of 1/2,2-1}\\
        \Res_{z_2=-\frac{k}{2}}\omega_{\frac12,2}(z_1,z_2)e^{z_2L_2}=&\mathfrak{b}L_2\frac{(-1)^k}{k}e^{-\frac{k}{2}L_2}\left(\frac{dz_1}{(z_1-\frac{k}{2})^2}+\frac{dz_1}{(z_1+\frac{k}{2})^2}\right),\label{inverse of 1/2,2-2}\\
        \Res_{z_2=-z_1}\omega_{\frac12,2}(z_1,z_2)e^{z_2L_2}=&-\mathfrak{b}\pi e^{-L_2 z_1}\left(\frac{L_2^2\,dz_1}{z_1\sin2\pi z_1}+\frac{L_2\,dz_1}{z_1^2\sin2\pi z_1}+\frac{2\pi L_2\cos2\pi z_1\,dz_1}{z_1\sin^22\pi z_1}\right).\label{inverse of 1/2,2-3}
    \end{align}
    %cancels that in \eqref{inverse of 1/2,2-3}
Note that the pole at $z=\frac{k}{2}$ in \eqref{inverse of 1/2,2-2} cancels the corresponding pole in \eqref{inverse of 1/2,2-3}; hence, $\mathcal{L}^{-1}_{L_2}.\;\omega_{\frac12,2}(z_1,z_2)$ has poles only at $z_1=0$ or $z_1=-\frac{k}{2}$.

Since we assumed $L_1\geq L_2$, one can now choose the contour for $\mathcal{L}^{-1}$ such that it encloses the origin and the negative real axis, regardless of the terms. From \eqref{inverse of 1/2,2-1} and \eqref{inverse of 1/2,2-2}, we have
\begin{align}
\Res_{z_1=0} \,e^{z_1L_1}\Res_{z_2=0}\omega_{\frac12,2}(z_1,z_2)\,e^{z_2L_2}=&2L_1L_2\mathfrak{b}\log\left(2\sinh\frac{\epsilon}{2}\right)-\frac{\mathfrak{b}L_1L_2^2}{2},\label{inverse of 1/2,2-4}\\
      \Res_{z_1=-\frac{k}{2}} \,e^{z_1L_1} \Res_{z_2=-\frac{k}{2}}\omega_{\frac12,2}(z_1,z_2)\,e^{z_2L_2}=&\mathfrak{b}L_1L_2\frac{(-1)^k}{k}e^{-\frac{k}{2}(L_1+L_2)}.\label{inverse of 1/2,2-5}
\end{align}
And for \eqref{inverse of 1/2,2-3}, we have contributions from $z_1=0$ and $z_1=-\frac{k}{2}$
    \begin{align*}
        \Res_{z_1=0}\,e^{z_1L_1} \Res_{z_2=-z_1}\omega_{\frac12,2}(z_1,z_2)e^{z_2L_2}=&-\frac{\mathfrak{b}}{2}L_1L_2(L_1-L_2),\\
        \Res_{z_1=-\frac{k}{2}}\,e^{z_1L_1} \Res_{z_2=-z_1}\omega_{\frac12,2}(z_1,z_2)e^{z_2L_2}=&\,\mathfrak{b}L_1L_2\frac{(-1)^k}{k}e^{-(L_1-L_2)\frac{k}{2}}.
    \end{align*}

Combining all these contributions, we obtain:
\begin{align*}
    (\hat{\mathcal{L}}_{L_1}^{-1} \hat{\mathcal{L}}_{L_2}^{-1}).\;\omega_{\frac12,2}=&\mathfrak{b}L_1L_2\left(2\log\left(2\sinh\frac{\epsilon}{2}\right)-\frac{L_1}{2}+\sum_{k\geq1}\frac{(-1)^k}{k}\left(e^{-(L_1+L_2)\frac{k}{2}}+e^{-(L_1-L_2)\frac{k}{2}}\right)\right)\nonumber\\
    =&\mathfrak{b}L_1L_2\left(2\log\left(2\sinh\frac{\epsilon}{2}\right)-\frac{L_1}{2}+\log\left(1+e^{-\frac{L_1+L_2}{2}}\right)-\log\left(1+e^{-\frac{L_1-L_2}{2}}\right)\right)\nonumber\\
    =&-\mathfrak{b}L_1L_2\left(\log\left(2\cosh\frac{L_1}{2}+2\cosh\frac{L_2}{2}\right)-2\log\left(2\sinh\frac{\epsilon}{2}\right)\right),
\end{align*}
where in the second equality we resummed the series in $k$ under the assumption that $L_1\geq L_2$. Thus, we arrive at the theorem by setting $b=1$, which means $\mathfrak{b}=-\frac{1}{\sqrt{2}}$. If we assume $L_1\leq L_2$, then we still obtain the same result by choosing a contour differently for the inverse Laplace transform of \eqref{inverse of 1/2,2-3}. Equivalently, one may repeat exactly the same argument by reversing the order of the inverse Laplace transforms, i.e.~applying $\mathcal{L}^{-1}_{L_1}\mathcal{L}^{-1}_{L_2}$, with an appropriate choice of contours. This explains how poles along the anti-diagonal should be treated.
\end{proof}

\end{document}
